\section{Conclusion}
\label{sec:conclusion}

Private compatibility does not have a sign-invariant effect on governments' incentives for formal harmonization. In the canonical three-country Cournot model, a regional member's pre-adoption payoff difference relative to international standardization is $\mathscr E-\mathscr D$. After the outsider adopts the bloc standard with residual marginal adaptation cost $d$, the difference becomes
\[
-\mathscr D+R(d,v).
\]
The equilibrium-derived residual rent $R(d,v)$ therefore separates three cases: actual reversal toward international standardization, stronger international-standardization incentives without reversal, and weaker international-standardization incentives.

The mechanism is not confined to complete cost elimination. Incomplete adaptation, selection-free outsider-only adoption, and an actual government ranking reversal coexist on a nonempty open set. But the political effect is not broadly portable. The initial regional advantage fails at the zero-network boundary, under one alternative singleton-network specification, and throughout the pre-specified differentiated-demand comparison under both Cournot and Bertrand competition, even though one-way adoption can remain feasible.

The coalition-stability application illustrates the institutional consequence without changing that scope. Under full bypass, an intermediate fixed-cost region makes international standardization uniquely stable in the canonical model. The symmetric high-cost regional stable set is more fragile because it depends on strict blocking and symmetry-induced indifference.

The paper's conclusion is therefore conditional. Private post-formation adaptation can alter the rents that sustain a formal standards coalition, but whether it pushes governments toward broader harmonization depends on how much rent the workaround removes and on the market structure generating those rents. The failure boundaries are part of the result, not exceptions to a general complementarity claim.
